\section{Methodology}
\label{sec:methodology}

\subsection{Data Sources}

Data is obtained via the FastF1 Python library, which provides access to the
official F1 timing database. We collect race results, qualifying times, lap-by-lap
telemetry (speed, throttle, brake, DRS activation), and weather station
measurements for all race weekends from 2020 through 2025—a total of approximately
120 race weekends and 2,400 driver-race records.

\subsection{Feature Engineering}

Features are organized into four groups:

\subsubsection{Lap Features (15 features)}
Derived from per-lap telemetry data: average lap time over the last 5 laps,
median and best lap time, standard deviation (consistency), pace trend (linear
slope), maximum and average speed, throttle/brake usage fractions, cornering and
straight-line speed, estimated fuel consumption rate, tire degradation rate,
and DRS activation count.

\subsubsection{Driver Features (20 features)}
Historical performance indicators: experience (years, total races), career win rate
and podium rate, points per race, DNF rate, rolling 5-race averages for finishing
position and points, grid-to-finish position delta, teammate comparison, best
circuit performance, and team stability indicators.

\subsubsection{Constructor Features (10 features)}
Team-level metrics: average and recent points, DNF rate, win count, reliability
score, engine speed ranking relative to competitor power units, and flags for
new-regulation seasons.

\subsubsection{Weather Features (8 features)}
Track and air temperature, humidity, wind speed, rainfall binary indicator, track
status code (0=dry to 3=very wet), temperature differential, and a composite
aggressiveness score that weights rain, wind, and humidity.

\subsection{Feature Selection}

A variance threshold ($\tau = 0.01$) removes near-constant features, followed by
correlation filtering that removes one member of each pair exceeding $r = 0.95$.
The top 30 features by variance are retained for modelling.

\subsection{Model Architectures}

\subsubsection{XGBoost}
An \texttt{XGBClassifier} with \texttt{multi:softprob} objective predicts
probability distributions over 20 position classes. Hyperparameters: max\_depth=6,
learning\_rate=0.1, n\_estimators=200, subsample=0.8, colsample\_bytree=0.8.
Early stopping with a 50-round patience is applied when validation data is available.

\subsubsection{Neural Network}
A feed-forward network with architecture [input, 256, 128, 64, 20]. Each hidden
layer uses ReLU activation, BatchNorm, and Dropout (0.3). Training uses Adam
optimizer (lr=0.001, weight\_decay=1e-4) with cross-entropy loss and early stopping
after 20 epochs without validation improvement.

\subsubsection{Baseline}
Predicts race position as grid position adjusted by the historical mean
grid-to-finish delta for each driver. Positions are clipped to [1, 20].

\subsection{Training and Evaluation Protocol}

Data is split temporally: 2020--2023 (train), 2024 (validation), 2025 (test).
This prevents any future information from influencing model fitting. Five-fold
temporal cross-validation uses an expanding window where each fold adds one
more season to the training set.

Evaluation metrics include mean absolute error (MAE), RMSE, exact position
accuracy, top-3 and top-5 accuracy (fraction of predictions within $k$ positions
of truth), Spearman rank correlation, and normalized discounted cumulative gain
(NDCG@10).
